\newcommand{\objone}{
\begin{itemize}
    \item \textbf{Input Handling:}  
    To write a Java program that takes input for an array of integers from the user.  
    The program should be able to handle any valid size of the array entered by the user.

    \item \textbf{Array Splitting:}  
    To divide the array into two halves based on its length.  
    If the array has an odd number of elements, the extra element belongs to the second half.

    \item \textbf{Ascending Sort:}  
    To sort the first half of the array in ascending order using Java’s built-in sorting method.  
    This ensures that the initial part of the array is organized from smallest to largest.

    \item \textbf{Descending Sort:}  
    To sort the second half of the array in descending order.  
    This requires converting the elements to an object array to use reverse order sorting.

    \item \textbf{Combining Halves:}  
    To combine both sorted halves back into the original array.  
    Care should be taken to maintain the order of each sorted half during merging.

    \item \textbf{Output Display:}  
    To display the final array with the first half sorted in ascending order  
    and the second half in descending order as the final output to the user.
\end{itemize}
}

\newcommand{\objtwo}{
\begin{itemize}
    \item \textbf{Matrix Input:}  
    To develop a Java program that reads a 2D matrix from the user.  
    The program should accept user-defined numbers of rows and columns.

    \item \textbf{Storing Data:}  
    To store the entered matrix elements into a 2D array structure.  
    This allows further processing for analysis and output.

    \item \textbf{Row Minimums:}  
    To compute the minimum element from each row of the matrix.  
    These values will be displayed alongside each respective row.

    \item \textbf{Column Minimums:}  
    To determine the smallest value from each column in the matrix.  
    These values will be printed separately after the matrix.

    \item \textbf{Formatted Output:}  
    To display the matrix elements clearly in a tabular form.  
    Row minimums are shown at the end of each row for readability.

    \item \textbf{Minimum Summary:}  
    To provide a summary of column minimums below the matrix.  
    This ensures the output reflects both row-wise and column-wise analysis.
\end{itemize}
}

\newcommand{\objthree}{
\begin{itemize}
    \item \textbf{User Input:}  
    To create a Java program that accepts a string from the user.  
    The input may contain alphabets, numbers, symbols, and spaces.

    \item \textbf{Character Analysis:}  
    To process each character in the string and identify whether it is a vowel, consonant, or non-letter.  
    The identification is case-insensitive.

    \item \textbf{Vowel Retention:}  
    To retain all vowels (A, E, I, O, U and their lowercase forms) in the output.  
    These characters are preserved in their original order.

    \item \textbf{Consonant Removal:}  
    To remove all consonant characters from the string.  
    Only alphabetic consonants are excluded; other characters remain unchanged.

    \item \textbf{Result Construction:}  
    To build the resulting string by appending valid characters (vowels and non-letters).  
    A `StringBuilder` or similar structure may be used for efficient concatenation.

    \item \textbf{Output Display:}  
    To display the final string after removing consonants.  
    The output reflects only vowels and non-alphabet characters from the original input.
\end{itemize}
}

\newcommand{\objfour}{
\begin{itemize}

    \item \textbf{Class Definition:} 
    Define a class named \texttt{MyPoint} to model a point in 2D space with integer coordinates \texttt{x} and \texttt{y}.

    \item \textbf{Default Constructor:} 
    Create a no-argument constructor to initialize the point at (0, 0), representing the origin.

    \item \textbf{Parameterized Constructor:} 
    Define an overloaded constructor that accepts two integers to set custom values of \texttt{x} and \texttt{y}.

    \item \textbf{Setting Values:} 
    Implement a method \texttt{setXY()} that assigns values to both \texttt{x} and \texttt{y} simultaneously.

    \item \textbf{Returning Coordinates:} 
    Implement \texttt{getXY()} to return a 2-element array containing the current values of \texttt{x} and \texttt{y}.

    \item \textbf{Custom String Output:} 
    Override the \texttt{toString()} method to return a string in the format \texttt{"(x, y)"}.

    \item \textbf{Distance Calculation:} 
    Write a method \texttt{distance(int x, int y)} that computes the distance from the current point to another point using the distance formula.

    \item \textbf{Test Class Creation:} 
    Define a separate class \texttt{TestMyPoint} containing the \texttt{main()} method to test the functionality of the \texttt{MyPoint} class.

    \item \textbf{Object Instantiation and Testing:} 
    In the test class, create one or more \texttt{MyPoint} objects and call all public methods to verify expected behavior.

    \item \textbf{Display Results:} 
    Print the output from all test cases, including coordinates, string representation, and distance values.

\end{itemize}
}

\newcommand{\objfive}{
\begin{itemize}

    \item \textbf{Define Superclass \texttt{Person}:}
    Create a superclass named \texttt{Person} with common attributes such as \texttt{name}, \texttt{address}, and methods for setting and displaying information.

    \item \textbf{Create Subclass \texttt{Student}:}
    Extend the \texttt{Person} class to form a \texttt{Student} class with additional attributes like \texttt{program}, \texttt{year}, and \texttt{fee}.

    \item \textbf{Create Subclass \texttt{Staff}:}
    Similarly, create another subclass \texttt{Staff} with specific attributes such as \texttt{school} and \texttt{pay}.

    \item \textbf{Implement Constructors:}
    Define constructors in all classes for proper initialization of attributes, using \texttt{super} keyword where needed to invoke the superclass constructor.

    \item \textbf{Override \texttt{toString()} Methods:}
    Provide meaningful string representations for each class by overriding the \texttt{toString()} method to include both inherited and specific attributes.

    \item \textbf{Use Setter and Getter Methods:}
    Implement getter and setter methods to allow controlled access to private data fields and to demonstrate encapsulation.

    \item \textbf{Instantiate Classes in Main:}
    In a separate driver class, create objects of \texttt{Person}, \texttt{Student}, and \texttt{Staff}, initializing them with suitable values.

    \item \textbf{Demonstrate Polymorphism:}
    Use polymorphism by referring to subclass objects using superclass references and calling overridden methods.

    \item \textbf{Display Object Details:}
    Call the \texttt{toString()} methods for each object and print their information on the console for verification.

    \item \textbf{Test and Validate:}
    Verify the correctness of constructor chaining, method overriding, and data handling through sample output and debugging if needed.

\end{itemize}
}

\newcommand{\objsix}{
\begin{itemize}
    \item \textbf{Understanding Inheritance:} 
    To demonstrate how a derived class can extend the functionality of a base class by inheriting its properties.

    \item \textbf{Constructor Usage:} 
    To learn how to initialize variables of both base and derived classes using constructors, including the use of the \texttt{super} keyword.

    \item \textbf{Method Overriding:} 
    To practice method overriding by redefining the \texttt{getVolume()} method in the derived class to provide specific functionality.

    \item \textbf{Applying Mathematical Concepts:} 
    To apply formulas for volume calculation of a cube and a cylinder using class methods.

    \item \textbf{Encapsulation and Reusability:} 
    To reinforce the importance of encapsulating data and designing reusable components using object-oriented principles.

    \item \textbf{Polymorphism Implementation:} 
    To understand how runtime polymorphism works through method overriding in an inheritance hierarchy.
\end{itemize}
}

\newcommand{\objseven}{
\begin{itemize}
    \item \textbf{Understanding Interfaces:} 
    To learn how to define and use interfaces in Java for specifying a contract of functionalities.

    \item \textbf{Abstraction in OOP:} 
    To implement abstraction by separating engine behavior definitions from their implementation using interfaces.

    \item \textbf{Implementing Interface Methods:} 
    To understand that any class implementing an interface must provide concrete implementations for all declared methods.

    \item \textbf{Encapsulation of Engine Properties:} 
    To encapsulate engine-related data such as \texttt{speed} and \texttt{gear} within a class.

    \item \textbf{Customizing Behavior:} 
    To allow specific behavior for speed and gear changes by implementing methods like 
    \\\texttt{speedUp(value)} and \texttt{changeGear(value)}.

    \item \textbf{Designing Flexible Systems:} 
    To design flexible and extensible systems where multiple classes can implement the same interface with different logic.
\end{itemize}
}

\newcommand{\objeight}{
\begin{itemize}

    \item \textbf{Understanding Exception Handling:} \\
    To learn how Java's exception mechanism works by using \texttt{try}, \texttt{catch}, and \texttt{finally} blocks. This helps manage runtime anomalies without abruptly terminating the program.

    \item \textbf{Handling Multiple Exceptions:} \\
    To demonstrate the handling of multiple types of exceptions (such as \texttt{ArrayIndexOutOfBoundsException},  using separate or multi-catch blocks in one program.

    \item \textbf{Improving Program Robustness:} \\
    To build fault-tolerant applications that anticipate and recover from common errors, resulting in more stable and user-friendly software.

    \item \textbf{Using Control Structures:} \\
    To integrate exception handling with conditional statements to define alternative execution paths and ensure uninterrupted program flow during error conditions.

    \item \textbf{Practicing Defensive Programming:} \\
    To follow best practices for error prevention and validation of user inputs, reducing the likelihood of unhandled exceptions and promoting reliable code execution.

    \item \textbf{Using Finally Block for Cleanup:} \\
    To ensure that important cleanup operations (like closing resources) are executed regardless of whether an exception occurs or not.

    \item \textbf{Enhancing Code Debugging and Maintenance:} \\
    To use exception messages and stack traces to quickly locate the source of errors, simplifying the debugging and maintenance process of the application.

\end{itemize}
}

\newcommand{\objnine}{

\begin{itemize}
    \item \textbf{Understanding JDBC Connectivity:} 
    To learn how to establish a connection between a Java application and a relational database using JDBC (Java Database Connectivity).

    \item \textbf{Data Insertion into Database:} 
    To insert user input data (Registration Number, Student Name, City, and Contact Number) into a database table using SQL INSERT statements.

    \item \textbf{Retrieving and Displaying Data:} 
    To fetch and display all records from the \texttt{Student\_info} table in a structured, tabular format using SQL SELECT queries.

    \item \textbf{Prepared Statements and Input Handling:} 
    To use \texttt{PreparedStatement} for secure and efficient execution of SQL commands while handling user input.

    \item \textbf{Database Design and Table Structure:} 
    To design a basic database schema that models student data in a tabular form.

    \item \textbf{Developing Practical Database Applications:} 
    To apply theoretical knowledge of databases in practical applications that involve both front-end input and back-end storage.
\end{itemize}
}

\newcommand{\objten}{
\begin{itemize}
    \item \textbf{Understanding Array Operations:} 
    To develop a clear understanding of how arrays work in Java, including traversal and manipulation.

    \item \textbf{Removing Elements Manually:} 
    To implement logic for removing a specific element from an array, as Java arrays have fixed size and no built-in remove method.

    \item \textbf{Using Conditional Statements:} 
    To practice decision-making using conditions for identifying the target element to remove.

    \item \textbf{Enhancing Logical Thinking:} 
    To improve algorithmic thinking by handling array reconstruction after removing an element.

    \item \textbf{Working with User Input:} 
    To read dynamic input from users for both the array and the element to be removed.
\end{itemize}
}

\newcommand{\objeleven}{
\begin{itemize}
    \item \textbf{Understanding Array Structure:}
    To gain a fundamental understanding of how arrays are organized and indexed in Java.

    \item \textbf{Element Insertion Logic:}
    To implement logic for inserting a new element into an existing array at a specific position by shifting elements.

    \item \textbf{Working with Fixed-Size Data Structures:}
    To handle fixed-size arrays in Java and create new arrays when modifications are required.

    \item \textbf{Input Validation and User Interaction:}
    To accept user input for the array, the element, and the target position, ensuring proper handling of input constraints.

    \item \textbf{Problem-Solving with Arrays:}
    To enhance problem-solving and algorithmic skills through manual manipulation of data structures.
\end{itemize}
}

\newcommand{\objtwelve}{
\begin{itemize}
    \item \textbf{Understanding Array Traversal:}
    To practice iterating over arrays using nested loops for pairwise comparison of elements.

    \item \textbf{Solving a Real-World Problem:}
    To identify element pairs that satisfy a given condition (sum = target), a common problem in algorithm and interview practice.

    \item \textbf{Improving Logical Thinking:}
    To enhance problem-solving and algorithmic thinking by implementing efficient search strategies.

    \item \textbf{Working with Conditional Statements:}
    To use conditional logic to check for pairs that satisfy a mathematical condition.

    \item \textbf{User Interaction and Input Handling:}
    To accept dynamic user input for the array elements and target sum, ensuring real-time processing of data.
\end{itemize}
}

\newcommand{\objthirteen}{
\begin{itemize}
    \item \textbf{Understanding Array Manipulation:}
    To learn how to manipulate array data by identifying and eliminating duplicate values.

    \item \textbf{Practicing Sorting Techniques:}
    To use array sorting as a strategy to simplify the detection and removal of duplicates.

    \item \textbf{Improving Algorithmic Thinking:}
    To implement logic that filters unique values and constructs a new array of non-repeating elements.

    \item \textbf{Calculating Effective Array Size:}
    To determine and return the updated length of the array after removing duplicates.

    \item \textbf{Handling User Input:}
    To accept dynamic input from users and process it effectively through iterative and conditional logic.
\end{itemize}
}

\newcommand{\objfourteen}{
\begin{itemize}
    \item \textbf{Understanding Sequence Detection:}
    To learn how to identify consecutive sequences within an unsorted array of integers.

    \item \textbf{Working with Unordered Data:}
    To handle and process unsorted data using data structures such as arrays and sets.

    \item \textbf{Applying Efficient Algorithms:}
    To implement an optimized algorithm (using HashSet or sorting) to find the longest sequence without unnecessary computation.

    \item \textbf{Enhancing Problem-Solving Skills:}
    To develop algorithmic thinking by solving a classic array-based problem commonly asked in coding interviews.

    \item \textbf{Practicing Java Collections:}
    To apply Java collections like \texttt{HashSet} for fast lookups and sequence tracking.

    \item \textbf{Handling Edge Cases:}
    To account for duplicate elements and disconnected number sequences during processing.
\end{itemize}
}

\newcommand{\objfifteen}{
\begin{itemize}
    \item \textbf{Understanding String Comparison:}
    To learn how to compare two strings lexicographically using built-in Java methods.

    \item \textbf{Exploring the \texttt{compareTo()} Method:}
    To demonstrate the use of the \texttt{String.compareTo()} method for determining the ordering of strings based on Unicode values.

    \item \textbf{Handling Alphabetical Ordering:}
    To understand how strings are ordered in dictionaries and how Java handles such ordering internally.

    \item \textbf{Practicing Conditional Statements:}
    To implement logic that interprets the result of a lexicographical comparison and outputs meaningful messages.

    \item \textbf{Strengthening String Manipulation Skills:}
    To build proficiency in working with Java's \texttt{String} class and related methods.
\end{itemize}
}

\newcommand{\objsixteen}{
\begin{itemize}
    \item \textbf{Understanding String Comparison Techniques:}
    To learn how to compare specific parts (substrings) of two strings rather than comparing entire strings.

    \item \textbf{Using \texttt{regionMatches()} Method:}
    To explore the usage of the \texttt{regionMatches()} method in the \texttt{String} class for partial string comparison.

    \item \textbf{Improving String Manipulation Skills:}
    To gain hands-on experience in working with string indexing and substring matching.

    \item \textbf{Practicing Conditional Logic:}
    To apply conditional checks based on the outcome of region comparison for logical decision making.

    \item \textbf{Enhancing Debugging and Output Interpretation:}
    To interpret substring match results and print them clearly in a formatted manner.
\end{itemize}
}

\newcommand{\objseventeen}{
\begin{itemize}
    \item \textbf{Understanding String Permutations with Repetition:}  
    To learn how to generate all possible string permutations when characters are allowed to repeat.

    \item \textbf{Implementing Nested Loops for Permutation Logic:}  
    To apply nested loops effectively for generating combinations of characters from the input string.

    \item \textbf{Practicing Character Access and Concatenation:}  
    To improve skills in accessing individual characters and forming new string patterns.

    \item \textbf{Improving Loop Control and Iteration:}  
    To gain hands-on experience in iterating through multiple levels to generate repeated patterns.

    \item \textbf{Enhancing Output Formatting Skills:}  
    To display generated permutations in a clear, structured, and readable format.
\end{itemize}
}

\newcommand{\objeighteen}{
\begin{itemize}
    \item \textbf{Understanding String Tokenization:}  
    To learn how to break a string into individual words using built-in methods like \texttt{split()}.

    \item \textbf{Handling Whitespace and Input Variations:}  
    To correctly manage extra spaces, tabs, and newline characters when counting words.

    \item \textbf{Practicing Conditional and Null Checks:}  
    To implement safe checks for null or empty strings to avoid runtime errors.

    \item \textbf{Improving User Interaction via Console Input:}  
    To use \texttt{Scanner} for accepting user input and processing it efficiently.

    \item \textbf{Enhancing Output Formatting and Validation:}  
    To display the total word count in a user-friendly and formatted manner.
\end{itemize}
}

\newcommand{\objnineteen}{
\begin{itemize}
    \item \textbf{Understanding Object as Constructor Argument:}  
    To learn how to use an object as a parameter in a constructor for initializing another object's data members.

    \item \textbf{Applying Class and Object Concepts:}  
    To reinforce knowledge of class creation, object instantiation, and constructor overloading.

    \item \textbf{Practicing Data Encapsulation:}  
    To encapsulate box dimensions and provide a method to compute volume based on internal data.

    \item \textbf{Enhancing Method Creation Skills:}  
    To define and invoke a method that performs a computation (volume) and returns the result.

    \item \textbf{Improving Output Presentation:}  
    To display the computed volume in a clear and formatted way from the \texttt{main()} method.
\end{itemize}
}


\newcommand{\objtwenty}{
\begin{itemize}
    \item \textbf{Understanding Constructor Overloading:}  
    To learn how to define and use both default and parameterized constructors, including object copy constructors.

    \item \textbf{Applying Object-Oriented Principles:}  
    To reinforce core OOP concepts like encapsulation, object passing, and method return types.

    \item \textbf{Practicing Salary Component Calculations:}  
    To implement a method that calculates PF (Provident Fund), allowances, and other salary components.

    \item \textbf{Returning Multiple Values via Objects:}  
    To return multiple calculated values using an object instead of primitive return types.

    \item \textbf{Enhancing Output Formatting and Data Handling:}  
    To display employee details and salary breakdown in a well-structured and readable format.
\end{itemize}
}

\newcommand{\objtwentyone}{
\begin{itemize}
    \item \textbf{Understanding User-Defined Exceptions:}  
    To learn how to create custom exception classes in Java by extending the \texttt{Exception} class.

    \item \textbf{Applying Exception Handling Mechanism:}  
    To implement \texttt{try-catch} blocks to detect and handle exceptional cases such as negative array sizes.

    \item \textbf{Improving Input Validation:}  
    To validate user inputs and throw exceptions proactively when inputs do not meet defined constraints.

    \item \textbf{Practicing Constructor Overriding in Exception Class:}  
    To define custom constructors and messages within user-defined exception classes.

    \item \textbf{Enhancing Robustness and User Feedback:}  
    To ensure the program fails gracefully with meaningful error messages instead of crashing on invalid input.
\end{itemize}
}

\newcommand{\objtwentytwo}{
\begin{itemize}
    \item \textbf{Understanding Object Comparison:}  
    To implement a method that compares two objects' data members to determine logical equality.

    \item \textbf{Practicing Constructor Initialization:}  
    To learn how to use parameterized constructors to initialize object attributes effectively.

    \item \textbf{Applying Boolean Logic and Return Types:}  
    To return a boolean value from a user-defined method based on conditional checks.

    \item \textbf{Reinforcing Object-Oriented Principles:}  
    To develop a clear understanding of object creation, passing objects as parameters, and encapsulating logic.

    \item \textbf{Displaying Logical Results Clearly:}  
    To display meaningful output in the \texttt{main()} method indicating whether two objects are equal or not.
\end{itemize}
}

\newcommand{\objtwentythree}{
\begin{itemize}
    \item \textbf{Understanding Abstract Classes and Methods:}  
    To learn how to define and use abstract classes and abstract methods in Java for enforcing method implementation in derived classes.

    \item \textbf{Practicing Inheritance and Method Overriding:}  
    To implement inheritance by creating subclasses that extend an abstract class and override its methods.

    \item \textbf{Implementing Polymorphic Behavior:}  
    To use dynamic method dispatch by invoking overridden methods using base class references.

    \item \textbf{Designing Role-Specific Salary Calculation:}  
    To define and implement different logic for calculating net salary based on employee type (Manager or Clerk).

    \item \textbf{Enhancing Data Encapsulation and Output Clarity:}  
    To encapsulate employee properties and ensure well-formatted display of employee details using overridden methods.
\end{itemize}
}